\SourcePage{212}{215}
\section{The Bohr--Sommerfeld quantization rule}

States belonging to the discrete energy spectrum become quasiclassical when
the quantum number $n$, which gives the ordinal number of the state, is large.
Indeed, this number specifies the number of nodes of the eigenfunction (see
§~21). The separation of two successive nodes, however, has the same order of
magnitude as the de Broglie wavelength. Hence, for large $n$ this separation
is small, and the wavelength is short in comparison with the dimensions of the
region of motion.

Let us derive the condition that fixes the quantum energy levels in the
quasiclassical case. Consider a particle executing bounded one-dimensional
motion in a potential well. The classically accessible interval
$b\leq x\leq a$ is bounded by two turning points\footnote{In classical
mechanics a particle in this field would move periodically, with period (the
time required to travel from $b$ to $a$ and back)
\[
  T=2\int_b^a\frac{dx}{v}
\]
($v$ is the particle velocity).}.

\SourcePage{213}{216}
By rule (47.5), the boundary condition at $x=b$ gives, in the region to its
right, the wave function
\begin{equation}
  \psi=\frac C{\sqrt p}\cos\left(\frac1\hbar\int_b^x p\,dx-\frac\pi4\right).
  \tag{48.1}
\end{equation}
Applying the same rule to the region on the left of $x=a$, we obtain another
form of the same wave function:
\[
  \psi=\frac{C'}{\sqrt p}\cos\left(\frac1\hbar\int_x^a p\,dx-\frac\pi4\right).
\]
For these expressions to agree throughout the interval, the sum of their
phases, which is constant, must be an integral multiple of $\pi$:
\[
  \frac1\hbar\int_b^a p\,dx-\frac\pi2=n\pi
\]
(with $C'=(-1)^nC$). It follows that
\begin{equation}
  \frac1{2\pi\hbar}\oint p\,dx=n+\frac12,
  \tag{48.2}
\end{equation}
where $\oint p\,dx=2\int_b^a p\,dx$ is evaluated over a complete period of
the particle's classical motion. This is the condition determining the
stationary states in the quasiclassical regime. It corresponds to the
Bohr--Sommerfeld quantization rule of the old quantum theory.

The quantity $I=(1/2\pi)\oint p\,dx$ is called the adiabatic invariant (see
I, §~49), and therefore (48.2) may be written
\[
  I(E)=\hbar(n+1/2).
\]
It was already noted in §~41 that, if the parameters change sufficiently
slowly, or ``adiabatically,'' a system remains in the same quantum state---in
the present case, the state having a particular $n$. We now see that in the
quasiclassical limit this statement becomes the classical theorem asserting
the constancy of the adiabatic invariant under a slow change of parameters.

\SourcePage{214}{217}
It is readily seen that the integer $n$ equals the number of zeros of the wave
function and is therefore the ordinal number of the stationary state. In
fact, the phase in (48.1) increases from $-\pi/4$ at $x=b$ to
$(n+1/4)\pi$ at $x=a$, so the cosine vanishes $n$ times within this interval
(outside the range $b\ll x\ll a$, the wave function decreases monotonically
and has no zeros at finite distances)\footnote{Strictly, the zeros must be
counted using the exact form of the wave function near the turning points. An
analysis of this kind confirms the stated result.}.

As explained above, $n$ is large in the quasiclassical case. Nevertheless,
retaining the term $1/2$ beside $n$ in (48.2) is legitimate. Inclusion of the
next correction terms in the phases of the wave functions would add to the
right-hand side of (48.2) only terms of order $\lambda/L$, which are small
compared with unity (see the remark at the end of §~46)\footnote{In certain instances, the exact formula for the energy levels
$E(n)$ (viewed as a function of the quantum number $n$), obtained from the
exact Schrödinger equation, retains the same form as $n\to\infty$.  Examples
are furnished by the energy levels in a Coulomb field and those of the
harmonic oscillator.  It is therefore natural that in these instances the
quantization rule (48.2), which applies for large $n$, yields for $E(n)$ a
formula identical to the exact one.}
.

To normalize the wave function, it is sufficient to integrate $|\psi|^2$
only over $b\leq x\leq a$, since $\psi(x)$ decays exponentially outside this
interval. The argument of the cosine in (48.1) varies rapidly, so to the
required accuracy its square may be replaced by its mean value, $1/2$. Thus
\[
  \int|\psi|^2dx\approx\frac{C^2}{2}\int_b^a\frac{dx}{p(x)}
  =\frac{\pi C^2}{2m\omega}=1,
\]
where $\omega=2\pi/T$ is the frequency of the classical periodic motion. The
normalized quasiclassical function is consequently
\begingroup\small
\begin{equation}
  \psi=2\sqrt{\frac{m\omega}{\pi p}}
  \cos\left(\frac1\hbar\int_b^x p\,dx-\frac\pi4\right).
  \tag{48.3}
\end{equation}
\endgroup
It should be kept in mind that $\omega$ depends on the energy and, in
general, differs from one level to another.

Relation (48.2) also admits another interpretation. The integral
$\oint p\,dx$ is the area enclosed by the particle's closed

\SourcePage{215}{218}
classical phase trajectory (that is, by the curve in the $p,x$ plane---the
particle's phase space). If this area is partitioned into cells, each of area
$2\pi\hbar$, their total number is $n$. But $n$ is the number of quantum
states whose energies do not exceed the specified value (the value
corresponding to the phase trajectory under consideration). Thus, in the
quasiclassical case, one may associate with every quantum state a phase-space
cell of area $2\pi\hbar$. Equivalently, the number of states assigned to a
phase-space area element $\Delta p\Delta x$ is
\begin{equation}
  \frac{\Delta p\Delta x}{2\pi\hbar}.
  \tag{48.4}
\end{equation}
If the wave vector $k=p/\hbar$ is used in place of the momentum, this number
takes the form $\Delta k\Delta x/2\pi$. As expected, it agrees with the known
expression for the number of normal modes of a wave field (see II, §~52).


The general pattern in which levels are distributed through the energy
spectrum can be deduced from the quantization rule (48.2). Let $\Delta E$
denote the separation of two adjacent levels, whose quantum numbers $n$ thus
differ by one. Since, for large $n$, $\Delta E$ is small in comparison with
the level energy itself, (48.2) gives
\[
  \Delta E\oint\frac{\partial p}{\partial E}\,dx=2\pi\hbar.
\]
But $\partial E/\partial p=v$, and hence
\[
  \oint\frac{\partial p}{\partial E}\,dx=\oint\frac{dx}{v}=T.
\]
We therefore find
\begin{equation}
  \Delta E=\frac{2\pi\hbar}{T}=\hbar\omega.
  \tag{48.5}
\end{equation}

Thus the spacing of adjacent levels is $\hbar\omega$. For a sequence of
nearby levels whose differences in $n$ are small compared with $n$ itself,
the corresponding frequencies $\omega$ may be regarded as approximately
equal. It follows that, within any small portion of the quasiclassical part of
the spectrum, the levels are equidistant, with the common interval
$\hbar\omega$.
This result is also foreseeable: in the quasiclassical regime, every frequency
associated with a transition between energy levels must be an integer multiple
of the classical frequency $\omega$.

\SourcePage{216}{219}

It is useful to trace the limiting behavior of the matrix elements of an
arbitrary physical quantity $f$. We start from the requirement that the mean
value of $f$ in a quantum state must tend simply to its classical value,
provided that, in the limit, the state itself describes motion of the particle
along a definite trajectory. Such a state is represented by a wave packet
(see §~6), formed by superposing stationary states with neighboring energies.
Its wave function has the form
\[
  \psi=\sum_n a_n\psi_n,
\]
where the coefficients $a_n$ differ appreciably from zero only within an
interval $\Delta n$ of quantum-number values such that
$1\ll\Delta n\ll n$. The numbers $n$ are assumed large, as required for the
stationary states to be quasiclassical. By definition, the mean value of $f$
is
\[
  \bar f=\int\psi^*f\psi\,dx
  =\sum_n\sum_m a_m^*a_nf_{mn}e^{i\omega_{mn}t},
\]
or, replacing the summation over $n,m$ by summations over $n$ and the
difference $s=m-n$, we obtain
\[
  \bar f=\sum_n\sum_s a_{n+s}^*a_nf_{n+s,n}e^{i\omega_st},
\]
where, in accordance with (48.5), $\omega_{mn}=\omega s$.

When evaluated with quasiclassical wave functions, the matrix elements
$f_{mn}$ decrease rapidly in magnitude as the difference $m-n$ grows, while
varying slowly with $n$ itself when $m-n$ is fixed. We may consequently write,
to the same approximation,
\[
  \bar f=\sum_n\sum_s a_n^*a_nf_se^{i\omega st}
  =\sum_n|a_n|^2\sum_sf_se^{i\omega st},
\]
where $f_s=f_{n+s,n}$ and $n$ denotes some mean value of the quantum number in
the interval $\Delta n$. Since $\sum_n|a_n|^2=1$, it follows that
\[
  \bar f=\sum_sf_se^{i\omega st}.
\]

\SourcePage{217}{220}
The resulting sum has the form of an ordinary Fourier series. Since $\bar f$
must approach the classical quantity $f(t)$ in the limit, we conclude that
the matrix elements $f_{mn}$ tend to the coefficients $f_{m-n}$ in the
Fourier-series expansion of the classical function $f(t)$.

Similarly, matrix elements for transitions between states in the continuous
spectrum tend to the components in the Fourier-integral expansion of $f(t)$.
Here the wave functions of the stationary states must be normalized to a
delta function of the energy divided by $\hbar$.

All the preceding results extend directly to systems with many degrees of
freedom undergoing finite motion for which the mechanical, or classical,
problem permits complete separation of variables in the Hamilton--Jacobi
method (the so-called conditionally periodic motion; see I, §~52).  After the
variables have been separated, the problem for each degree of freedom reduces
to a one-dimensional one, and the corresponding quantization conditions are
\begin{equation}
  \oint p_i\,dq_i=2\pi\hbar(n_i+\gamma_i),
  \tag{48.6}
\end{equation}
where the integral is taken over one period of the generalized coordinate
$q_i$, while $\gamma_i$ is a number of order unity determined by the form of
the boundary conditions for that degree of freedom.\footnote{For motion in a
central field, for example,
\[
  \oint p_r\,dr=2\pi\hbar(n_r+1/2),\qquad
  \oint p_\theta\,d\theta=2\pi\hbar(l-m+1/2),\qquad
  \oint p_\varphi\,d\varphi=2\pi\hbar m
\]
(where $n_r=n-l-1$ is the radial quantum number).  The last equality follows
simply because $p_\varphi$ is the $z$ component of the angular momentum and is
equal to $\hbar m$.}


For arbitrary multidimensional motion that is not conditionally periodic, a
formulation of the quasiclassical quantization conditions requires a more
subtle argument\footnote{See $J.,B.~Keller$ // Ann. Phys. 1958. V.~4.
P.~180.}. The notion of phase-space ``cells,'' however, always applies in the
same manner within the quasiclassical approximation. This is clear from the
connection, noted above, between this notion and the number of normal modes of
a wave field in a given spatial volume. In the general case of a system with
$s$ degrees of freedom, a phase-space volume element

\SourcePage{218}{221}
contains
\begin{equation}
  \Delta N=\frac{\Delta q_1\cdots\Delta q_s\,
  \Delta p_1\cdots\Delta p_s}{(2\pi\hbar)^s}
  \tag{48.7}
\end{equation}
quantum states\footnote{In particular, $d^3p/(2\pi\hbar)^3$ for a single
particle is the number of states associated with a momentum interval $d^3p$
in a unit volume of space. This accounts for the agreement between the two
interpretations of the plane-wave normalization (15.8) noted in the footnote
on p.~68.}.

\ProblemHeading

\textbf{1.} Find, approximately, the number of discrete energy levels for a
particle moving in a field $U(r)$ that satisfies the condition for the
quasiclassical approximation.


\SolutionHeading The number of states ``belonging'' to the phase-space volume
for momenta in the range $0\leq p\leq p_{\max}$ and particle coordinates in
the volume element $dV$ is
\[
  \frac{(4\pi/3)p_{\max}^3dV}{(2\pi\hbar)^3}.
\]
At a given $r$, the particle may have, in its classical motion, any momentum
that obeys $E=(p^2/2m)+U(r)\leq0$. Substitution of
$p_{\max}=\sqrt{-2mU(r)}$ gives the total number of discrete-spectrum states:
\[
  \frac{\sqrt2m^{3/2}}{3\pi^2\hbar^3}\int(-U)^{3/2}\,dV,
\]
where the integration extends over the part of space in which $U<0$. This
integral diverges, and the number of levels is infinite, if $U$ decreases at
infinity as $r^{-s}$ with $s<2$, in agreement with the results of §~18.

\textbf{2.} Find the corresponding number for a quasiclassical central field
$U(r)$ (V.~L.~Pokrovskii).

\SolutionHeading In a centrally symmetric field, the number of states differs
from the number of energy levels because the latter are degenerate with respect
to the orientation of the angular momentum. The required number may be found by
noting that the number of levels for a specified value of the angular momentum
$M$ is the same as the number of nondegenerate levels for one-dimensional
motion with potential energy
$U_{\text{eff}}=U(r)+M^2/(2mr^2)$. At a given $r$, the largest possible value of
the momentum $p_r$ for energies $E\leq0$ is
$p_{r\max}=\sqrt{-2mU_{\text{eff}}}$. Consequently, the number of states (that
is, the number of levels) is
\[
  \frac1{2\pi\hbar}\oint dr\,dp_r
  =\frac{\sqrt{2m}}{\pi\hbar}\int
  \sqrt{-U-\frac{M^2}{2mr^2}}\,dr.
\]
The required total number of discrete levels is obtained from this result by
integration with the measure $dM/\hbar$ (which, in the semiclassical case,
replaces summation over $l$), and is
\[
  \frac{\sqrt{2m}}{4\hbar^2}\int(-U)r\,dr.
\]

